\section{Introduction}
\label{sec:flex_motivation}

The preceding chapters addressed how agents can detect, evaluate, and recover from novelty in open-world environments. \Cref{ch:benchmarks} introduced benchmark environments and metrics for studying agent behavior under mismatch, and \cref{ch:rapid_learn} developed hybrid planning--learning methods that exploit action abstractions to achieve efficient adaptation when execution fails. Both contributions operate within the framework established in \cref{ch:foundations}: novelty arises from mismatch between the agent's internal models and the environment, and structured abstractions determine whether this mismatch can be managed effectively.

This chapter turns to the second dimension of the dissertation's central claim: structured abstractions over interaction can enable learned behaviors to generalize across objects, physical parameters, and robotic embodiments. While \cref{ch:rapid_learn} showed that action abstractions support efficient adaptation after failure, the work in this chapter asks whether interaction abstractions can be designed so that learned skills remain valid before failure occurs---that is, so that a policy trained on one object or robot transfers to new objects and robots without retraining.

The challenge is rooted in how robotic manipulation policies represent their interaction with the world. Many robot-learning approaches define policies over robot-centric representations: joint angles, end-effector poses, or motor torques. These representations tightly couple the learned behavior to a specific embodiment and control interface. A policy trained to open a drawer with one robot may encode not just the task-relevant forces and motions, but also that robot's kinematics, joint limits, and control dynamics. The result is a policy that solves the task on the training platform but fails when the robot changes, the object changes, or the physical parameters change, even when the underlying task physics remain the same.

FLEX addresses this limitation through object-centric force-space representations for articulated object manipulation. Policies are trained entirely in simulation without modeling any robot, using only the object's joint configuration and the applied force as state and action. An interactive joint parameter estimation procedure enables deployment on physical robots, and an operational space controller translates learned forces into robot commands. By representing manipulation at the level of forces applied to the object, FLEX captures physical invariants of the task while remaining agnostic to the robot that executes them.
